\section{Conclusion}

We define \emph{relaxed microarchitectural isolation} (RMI), a new security property allowing programs to share memory while restricting information leakage to that of non-speculative execution.
To enforce this type of secure speculation, we propose that processors use microarchitectural isolation to restrict the attacker's observations, and a combination of simple mechanisms for controlled speculation and program analysis.
We demonstrate our approach by building an end-to-end FPGA-based prototype running enclave programs with low performance overheads on a multicore out-of-order processor.
Our final platform, Citadel, is a novel design point at the intersection of TEEs and TES defense mechanisms, and a compelling trade-off between performance, usability, and security.
